\section*{章末练习}
\begingroup\small
\subsection*{口述概念}
\begin{enumerate}[nosep,leftmargin=*]
  \item 区分定义与命题。为什么“组合收益定义为 $R_p=\sum_iw_ir_i$”不能用样本数据证明，而凸组合界需要证明？
  \item 在 $A\Rightarrow B$ 中，分别说明必要条件与充分条件，并解释为什么通过一个必要检查不能推出研究结论可信。
  \item 什么是合法反例？为什么 $w=(2,0,0)$ 不是“删除权重非负性后凸组合界失效”的受控反例？
\end{enumerate}

\subsection*{笔试推导}
\begin{enumerate}[nosep,leftmargin=*]
  \item 写出全称命题 $\forall x\in\mathcal D,A(x)\Rightarrow B(x)$ 的否定，并据此说明一个反例必须满足哪两件事。
  \item 不看正文，证明：若 $w_i\geq0$ 且 $\sum_iw_i=1$，则 $\min_i r_i\leq\sum_iw_ir_i\leq\max_i r_i$；同时给出上下界取等条件。
  \item 在 $\sum_iw_i=1$ 下证明：若凸组合界对所有 $r\in\mathbb R^n$ 成立，则每个 $w_i\geq0$。指出证明中哪里使用了“对所有”。
\end{enumerate}

\subsection*{数值编程}
\begin{enumerate}[leftmargin=*]
  \item 计算 $r=(0.02,-0.01,0.03)$、$w=(0.5,0.3,0.2)$ 的逐资产贡献，再把权重改为 $(1.5,-0.5,0)$。
  \item 画出 $2w-1$、$1-2w$ 及二者最小值，求 $0\leq w\leq1$ 上最坏收益的最大值。
  \item 两期方差都为 4，分别取协方差 $1,0,-1$，比较和的标准差。
\end{enumerate}

\subsection*{研究判断}
\begin{enumerate}[nosep,leftmargin=*]
  \item 研究员说“回测包含手续费，因此结果可信”。把这句话写成逻辑箭头，判断手续费检查是必要、充分还是充要条件，并列出两个仍可能失败的原因。
  \item 报告把 100 美元 P\&L、2\% 日收益和 15\% 年化波动率相加得到“综合得分”。逐项审查量纲和时间尺度，给出一个有意义的替代报告结构。
  \item 某人用一百万组随机收益都没有找到凸组合界反例，于是声称定理已被程序证明。说明有限实验为何不足以证明全称命题；给出一般证明，再构造删除非负权重条件后的反例。
\end{enumerate}
\endgroup

\subsection*{基础衔接：独立计算}
两期收益方差都为 $4$，协方差为 $1$。求收益之和的方差和标准差；若误用独立缩放，得到多少？
